\chapter{Einstein Summation \& Bulk Array Operations}
\label{chap:einsum_bulk_ops}

\section{Einstein Summation Engine (\texttt{EINSUM})}
Sisal-2026 includes a built-in Einstein summation contraction engine. The syntax parses standard multi-tensor subscript notation:

\begin{lstlisting}[caption={Einstein Summation Syntax}]
C := EINSUM("ij,jk->ik", MatrixA, MatrixB)
\end{lstlisting}

\subsection{Lowering Fast-Paths}
The compiler inspects \texttt{EINSUM} subscript strings during AST lowering (\texttt{einsum\_lower.ml}) and routes contractions to hardware BLAS fast-paths:
\begin{itemize}
    \item \texttt{"ij,jk->ik"}: Matrix-Matrix Multiplication $\rightarrow$ \texttt{cblas\_dgemm} / \texttt{cblas\_sgemm}.
    \item \texttt{"ij,j->i"}: Matrix-Vector Multiplication $\rightarrow$ \texttt{cblas\_dgemv} / \texttt{cblas\_sgemv}.
    \item \texttt{"i,i->"}: Vector Dot Product $\rightarrow$ \texttt{cblas\_ddot} / \texttt{cblas\_sdot}.
    \item \texttt{"ij->ji"}: Matrix Transposition $\rightarrow$ $O(1)$ stride swap in dope vector metadata.
\end{itemize}

For general multi-index contractions not matching single 2D BLAS signatures, the compiler lowers the expression to an IF1 \texttt{EINSUM\_NODE} multi-loop iteration graph.

\subsection{High-Rank Tensors \& Contraction Path Optimization}
While \texttt{MATMUL} and \texttt{INNERPRODUCT} provide direct, concise syntax for standard 2D matrix operations, \texttt{EINSUM} serves two critical structural roles in tensor compilation:
\begin{enumerate}
    \item \textbf{High-Rank Tensor Contractions}: For 3D, 4D, and 5D attention graphs (e.g. $Q, K, V \in \mathbb{R}^{B \times H \times S \times D}$ in Grouped Query Attention), \texttt{EINSUM("bhqd,bhkd->bhqk", Q, K)} declaratively specifies contraction axes without manual reshape or permutation boilerplate.
    \item \textbf{Multi-Contraction Chain Optimization (Compiler Roadmap Target)}: When chaining multiple matrix/tensor products in sequence (e.g., $P \cdot V \cdot W_O$), exposing the contraction topology through \texttt{EINSUM} provides the structural foundation for \textbf{Dynamic Programming Contraction Path Optimization} (on the compiler optimizer roadmap). This pass will automatically evaluate intermediate tensor dimension products to select the optimal parenthesization order that minimizes total FLOP complexity (reducing $O(N^4)$ naive contractions down to optimal $O(N^3)$ sequences) before emitting BLAS kernel calls.
\end{enumerate}

\subsection{Tiled EINSUM Subscripts \& Unified Permutation (Wishlist Item)}
An upcoming design proposal on the Sisal-2026 compiler wishlist extends the \texttt{EINSUM} subscript string to natively unify \textbf{tensor contraction math}, \textbf{L1/L2 cache tiling}, and \textbf{loop permutation} into a single bracketed notation, eliminating separate imperative transformation scripts.

In this syntax, bracket nesting depth represents tile/cache depth (e.g. \texttt{(i 32)} represents a 32-element L2 block tile), while the left-to-right order of terms on the output side (\texttt{->}) dictates the execution loop nesting:

\begin{lstlisting}[language=Sisal,caption={Tiled EINSUM Subscript with Middle $k$ Loop}]
% Tiled Matrix Multiplication (32x64x16) with k as Middle Loop
C := EINSUM("((i 32) (j 64)) ((j 64) (k 16)) -> (i 32) (k 16) (j 64)", A, B)
\end{lstlisting}

\subsubsection{Step-by-Step Lowering Walkthrough (Pseudo-Fortran)}

To understand why placing $k$ as the middle loop delivers superior performance, consider the step-by-step compiler lowerings from naive untiled execution down to tiled $k$-in-the-middle micro-kernels:

\begin{enumerate}
    \item \textbf{Naive Untiled Lowering (\texttt{"ik, kj -> ij"} $\rightarrow$ $i \rightarrow j \rightarrow k$ order)}:
\begin{lstlisting}[language=Fortran]
! Naive Lowering: Reduction k in innermost loop (Strided access on B)
DO i = 1, M
   DO j = 1, N
      sum_val = 0.0d0
      DO k = 1, K
         sum_val = sum_val + A(i, k) * B(k, j)   ! B(k, j) jumps N elements every step
      END DO
      C(i, j) = sum_val
   END DO
END DO
\end{lstlisting}

    \item \textbf{Interchanged Lowering ($k$ as Middle Loop $\rightarrow$ $i \rightarrow k \rightarrow j$ order)}:
\begin{lstlisting}[language=Fortran]
! Interchanged: k in middle position (Scalar broadcast of A, unit-stride for B and C)
DO i = 1, M
   DO k = 1, K                       ! k is MIDDLE loop!
      a_val = A(i, k)                ! Broadcast: Load once into register
      DO j = 1, N                    ! Innermost loop over j (Vector unit-stride)
         C(i, j) = C(i, j) + a_val * B(k, j)
      END DO
   END DO
END DO
\end{lstlisting}

    \item \textbf{Tiled Lowering with $k$ in Middle (\texttt{"((i 32) (j 64)) ((j 64) (k 16)) -> (i 32) (k 16) (j 64)"})}:
\begin{lstlisting}[language=Fortran]
! Tiled & Interchanged: i_tile -> k_tile -> j_tile -> i_elem -> k_elem -> j_elem
DO i_tile = 1, M, 32
   DO k_tile = 1, K, 16                 ! k_tile is MIDDLE tile loop!
      DO j_tile = 1, N, 64
         
         ! Inner Micro-kernel (L1 Cache / SIMD Vector Registers)
         DO i_elem = i_tile, MIN(i_tile + 31, M)
            DO k_elem = k_tile, MIN(k_tile + 15, K)   ! k_elem is MIDDLE element loop!
               a_val = A(i_elem, k_elem)
               
               !DIR$ SIMD (Vectorized unit-stride loop over j_elem)
               DO j_elem = j_tile, MIN(j_tile + 63, N)
                  C(i_elem, j_elem) = C(i_elem, j_elem) + a_val * B(k_elem, j_elem)
               END DO
            END DO
         END DO

      END DO
   END DO
END DO
\end{lstlisting}

    \item \textbf{3D Batched MatMul \& 4D Attention Lowering (\texttt{"b ((m 32) (k 16)) (b ((k 16) (n 64))) -> b (m 32) (k 16) (n 64)"})}:
\begin{lstlisting}[language=Fortran]
! 3D Batched MatMul with Parallel Batch (b) & Middle k Loop
!$OMP PARALLEL DO COLLAPSE(2) PRIVATE(m_elem, k_elem, n_elem, a_val)
DO b = 1, B
   DO m_tile = 1, M, 32
      DO k_tile = 1, K, 16                 ! k_tile is MIDDLE tile loop!
         DO n_tile = 1, N, 64
            
            ! Inner Micro-kernel (Fits 100% in L1 Cache)
            DO m_elem = m_tile, MIN(m_tile + 31, M)
               DO k_elem = k_tile, MIN(k_tile + 15, K)   ! k_elem is MIDDLE element loop!
                  a_val = A(b, m_elem, k_elem)
                  
                  !DIR$ SIMD
                  DO n_elem = n_tile, MIN(n_tile + 63, N)
                     C(b, m_elem, n_elem) = C(b, m_elem, n_elem) + a_val * B(b, k_elem, n_elem)
                  END DO
               END DO
            END DO

         END DO
      END DO
   END DO
END DO
!$OMP END PARALLEL DO
\end{lstlisting}
    For 4D Transformer Attention ($Q \cdot K^T$ with Batch $b$, Heads $h$, Sequence $s$, Head Dim $d$), writing \texttt{Scores := EINSUM("b h ((s 64) (d 16)) (b h ((k 64) (d 16))) -> b h (s 64) (d 16) (k 64)", Q, K)} lowers to thread-parallel batch execution over $(b, h)$, outer tiles over Query sequence $(s)$, \textbf{head dimension $d$ reduced in middle registers}, and Key sequence $(k)$ streamed with unit-stride vector memory access.
\end{enumerate}

\subsubsection{Related Work \& Novelty Analysis}

A comparison of the Sisal-2026 tiled \texttt{EINSUM} proposal with state-of-the-art tensor compilers and programming frameworks highlights its unique design positioning:

\begin{itemize}
    \item \textbf{NumPy / PyTorch \texttt{einsum}}: Uses standard Einstein notation (\texttt{"ik,kj->ij"}). Expresses contraction math declaratively, but has \textbf{zero syntax for loop tiling, cache block sizes, or execution loop ordering}.
    \item \textbf{Halide (Ragan-Kelley et al., PLDI 2013) \& TVM (Chen et al., OSDI 2018)}: Pioneered the decoupling of computation algorithms from execution schedules. However, scheduling is imperative and separate from the math expression (e.g. \texttt{s.tile(i, i\_out, i\_in, 32).reorder(i\_out, k\_out, j\_out)} in an external Python script).
    \item \textbf{Einops (Rozenberg, ICLR 2022)}: Provides index rearrangement syntax (\texttt{"b (h w) -> b h w"}), but focuses exclusively on spatial array reshaping and tensor broadcasting---\textbf{cannot specify contraction reductions, cache block sizes, or SIMD loop order}.
    \item \textbf{Tensor Comprehensions (Vasilache et al., FAIR 2018) \& TACO (Kjolstad et al., OOPSLA 2017)}: Uses polyhedral compiler infrastructure (ISL) to discover tile sizes automatically during compilation. No inline syntax exists for the developer to specify tile bounds directly in the string.
    \item \textbf{Sisal-2026 Novelty}:
    \begin{enumerate}
        \item \textbf{Purely Functional Inline Scheduling}: Unifies math definition and hardware loop schedule in a single functional string expression.
        \item \textbf{Zero-Keyword Positional Ordering}: The position of bracketed terms on the RHS of \texttt{->} dictates outer/inner loop nesting and loop interchange (e.g. \texttt{-> (i 32) (k 16) (j 64)} places $k$ in the middle).
        \item \textbf{Cache Depth Mapping}: Bracket depth corresponds 1-to-1 to physical memory levels: \texttt{i} (flat DRAM), \texttt{(i 32)} (L2 Cache), \texttt{((i 128) 8)} (L1 / Vector Register).
    \end{enumerate}
\end{itemize}

\subsection{Higher-Order Epilogue \& Pre-Processing Fusion in EINSUM (Wishlist Item)}
Activation functions and transformations (such as \texttt{SwiGLU}, \texttt{Swish}, \texttt{GELU}, \texttt{ReLU}, \texttt{SiLU}, \texttt{RMSNorm}, \texttt{RoPE}, \texttt{BiasAdd}) can be captured \textbf{directly inside the subscript string} as named tokens without needing to pass them as extra function arguments to \texttt{EINSUM}:

\begin{lstlisting}[language=Sisal,caption={Fused SwiGLU FFN via EINSUM in Sisal-2026}]
function FusedSwiGLUFFN( X, W1, W2, W3 : array_dv[double] returns array_dv[double] )
  let
    % Activation token SwiGLU captured inside the string -- ZERO extra function args!
    Hidden := EINSUM(
      "((s 32) (h 64)) ((h 64) (d 16)), ((s 32) (h 64)) ((h 64) (d 16)) -> (s 32) (d 16) ((h 64) SwiGLU)",
      X, W1, X, W2
    )
  in
    EINSUM("((s 32) (d 16)) ((d 16) (h 64)) -> (s 32) (h 64) (d 16)", Hidden, W3)
  end let
end function
\end{lstlisting}

The subscript parser (\texttt{einsum\_lower.ml}) resolves built-in activation tokens or in-scope function identifiers directly from the AST symbol table. This unifies (1) tensor contraction math, (2) L1/L2 cache tiling, (3) loop interchange, (4) pre-processing load transformations, and (5) post-processing store activations \textbf{all inside a single self-contained \texttt{EINSUM} expression}!

\begin{note}[Work-in-Progress Disclaimer]
This documentation and manual are works in progress. Because compiler optimization passes, AST lowerings, and backend code generators are actively evolving, developers should examine the execution logs, generated C++ files, and benchmark results in \texttt{test/e2e/} to inspect the authoritative current state of the compiler's optimization capabilities.
\end{note}


\section{APL-Style Bulk Operations Battery (58 Primitives)}
Sisal-2026 provides 58 whole-array bulk primitives, eliminating manual parallel \texttt{forall} boilerplate code for standard array manipulations:

\begin{table}[h!]
\centering
\small
\begin{tabularx}{\textwidth}{l >{\raggedright\arraybackslash}p{4.0cm} >{\raggedright\arraybackslash}X}
\toprule
\textbf{Category} & \textbf{Primitive Names} & \textbf{Semantics \& Description} \\
\midrule
Reductions & \texttt{SUM}, \texttt{PRODUCT}, \texttt{LEAST}, \texttt{GREATEST}, \texttt{MEAN}, \texttt{VARIANCE}, \texttt{STDDEV}, \texttt{NORM} & Reduce array along default or specified axis into scalar value. \\
\addlinespace
Scans & \texttt{SCAN}, \texttt{CUMSUM}, \texttt{CUMPROD} & Prefix scan operation along leading axis. \\
\addlinespace
Structural & \texttt{TAKE}, \texttt{DROP}, \texttt{ROTATE}, \texttt{REVERSE}, \texttt{COMPRESS}, \texttt{RAVEL}, \texttt{EXPAND}, \texttt{SQUEEZE}, \texttt{PERMUTE} & Structural reshaping and axis permuting operations. \\
\addlinespace
Selection & \texttt{SORT}, \texttt{ARGMAX}, \texttt{ARGMIN}, \texttt{GRADE\_UP}, \texttt{GRADE\_DOWN}, \texttt{NONZERO}, \texttt{WHERE} & Sorting and index mask extraction. \\
\addlinespace
Higher-Order & \texttt{MAP}, \texttt{FOLDL}, \texttt{FOLDR}, \texttt{STENCIL}, \texttt{INNERPRODUCT}, \texttt{OUTERPRODUCT} & Functional combinators across multidimensional dope vectors. \\
\bottomrule
\end{tabularx}
\caption{APL Bulk Operation Battery Summary.}
\label{tab:bulk_ops}
\end{table}

\begin{lstlisting}[caption={Bulk Operation Battery Code Examples}]
let
  A := array_dv [1: 5.0, 2.0, 9.0, 1.0, 7.0];
  
  % Calculate Mean and Standard Deviation
  m := MEAN(A);
  s := STDDEV(A);
  
  % Sort and Grade Indices
  SortedA := SORT(A);          % [1.0, 2.0, 5.0, 7.0, 9.0]
  MaxIdx  := ARGMAX(A);        % 3 (index of value 9.0)
  
  % Stencil Convolution
  Smooth := STENCIL(A, 3, function(window : array_dv[real] returns real)
    (window[1] + window[2] + window[3]) / 3.0
  end function)
in
  SortedA, MaxIdx, Smooth
end let
\end{lstlisting}
